% IEEE TCAD -- Related Work section
% Included by main.tex; do not compile standalone.
%
% NOTE: as in introduction.tex, the citation keys below are bound to entries in
% refs.bib whose volume/page/date fields still need checking against the
% publisher record. No numeric claim about another author's work is quoted
% here; comparisons to prior measurements are deliberately qualitative, because
% this paper's own flow (Section~\ref{sec:impl}) is not the flow those papers
% used and the absolute values are not comparable.

\section{Related Work}
\label{sec:related}

Our experiment connects four main areas of prior research: permutation-based modes that separate the outer framework from the inner engine, making this kind of round substitution possible; ARX-based designs that provide our replacement round along with their well-known hardware and software trade-offs; existing benchmark studies of the NIST Lightweight Cryptography (LWC) candidates, which our eleven-design evaluation expands upon rather than repeats; and the published implementation literature on Ascon itself.

\subsection{Permutation-Based Modes and the Mode/Primitive Split}
The sponge construction \cite{sponge} and its duplex variant \cite{duplex} are built in two parts: the outer structure (the mode) and the inner mixing engine (the $b$-bit permutation). The mode governs how data is absorbed into the rate ($r$), while the capacity ($c$) is never touched by data directly---the property on which its security rests---and assumes the inner permutation is indistinguishable from a random one. Jovanovic \emph{et al.} \cite{jlm} proved that sponge-based authenticated encryption can achieve security beyond the birthday bound in the capacity, which forms the mathematical basis for Ascon's security claims \cite{caesar}.

Because the mode treats the permutation as a black box, the permutation can be exchanged without breaking the mode's mechanics. 
so Our design replaces Ascon's round function in full---both the 5-bit substitution layer and the linear diffusion layer---with SIPROUND from SIPHASH \cite{siphash}, retains Ascon's round-constant addition ahead of it, and narrows the state from five 64-bit words to four. The two layers are replaced together rather than separately because a single SIPROUND supplies, through its Addition--Rotation--XOR (ARX) operations, both the non-linearity and the inter-word mixing that Ascon divides between them.

Mixing modes and permutations is standard practice. Xoodyak \cite{xoodyak} runs its Cyclist mode over the Xoodoo permutation, while ISAP \cite{isap}, itself one of the nine finalists measured here, uses Ascon's permutation inside a different, leakage-resilient mode. However, this paper does the reverse: we hold a standardized mode completely fixed, swap its round function for one from a different design family, and measure the resulting costs across hardware and software targets.

\subsection{ARX Primitives and Their Cost Asymmetry}
Add-Rotate-XOR (ARX) designs, such as Salsa20 and ChaCha \cite{salsa,chacha}, BLAKE2 \cite{blake2}, SPECK \cite{simonspeck}, and SipHash \cite{siphash}, rely entirely on modular addition for nonlinearity. Because modern CPUs natively handle these operations without using lookup tables, ARX ciphers are extremely fast, compact, and naturally immune to the cache-timing side-channel attacks demonstrated by Bernstein against AES \cite{aes-cache}.

However, ARX designs struggle in hardware. Modular addition requires a carry propagation, which creates physical delays in the circuit. The SIMON and SPECK families \cite{simonspeck} clearly illustrate this divide: SIMON was built without addition for hardware efficiency, while SPECK relies on addition for software speed. SPARKLE, the only ARX-based NIST LWC finalist \cite{sparkle}, is built on a custom ARX-box called Alzette \cite{alzette}, designed to obtain strong cryptographic properties at bounded circuit depth. SPARKLE is the natural reference point for ARX in this field, though not a like-for-like one: its permutation is wider and performs considerably more work per invocation than SIPROUND over a 256-bit state, so Section~\ref{sec:results} reports it as context rather than as a matched ARX baseline.

The specific cost this paper measures---the 64-bit modular addition setting the critical path of a SipHash-derived round in silicon---has been identified independently. HAxSH \cite{haxsh} takes SipHash's round as given and attacks exactly that bottleneck, substituting approximate adders with shortened carry propagation for the exact ones and reporting substantial latency and power reductions from ASIC synthesis, with the resulting loss of exactness assessed through statistical testing and MILP-based differential and linear analysis. That work and this one share a diagnosis and differ in response: HAxSH changes the adder and keeps the construction exact nowhere, whereas this paper leaves SIPROUND untouched and asks what an unmodified ARX round costs when carried into a standardized AEAD mode. Its results are the natural point of comparison for any attempt to reduce the critical path reported in Section~\ref{sec:discussion}.

For our experiment, we substitute Ascon's s-box and linear diffusion layer with SIPROUND from SipHash \cite{siphash}. SipHash was originally designed as a short-input pseudorandom function for hash tables rather than for bulk encryption; the grounds for carrying its round into an AEAD permutation are set out above.

\subsection{Benchmarking of the NIST LWC Candidates}
The NIST Lightweight Cryptography (LWC) process produced extensive benchmarking data. Hardware comparisons were standardized using the GMU Hardware API \cite{lwc-api}---the interface the nine finalist cores measured here implement---allowing for consistent FPGA \cite{mohajerani} and ASIC \cite{aagaard} evaluations as summarized by NIST \cite{nist-lwc-report, nist-lwc-final}. In software, candidates were benchmarked across various microcontrollers \cite{renner,weatherley}. 

Our evaluation differs from these campaigns in two ways. First, our goal is not to rank candidates; instead, we use the eleven-design field as a calibration scale to measure the precise impact of our ARX substitution. Second, our software tests use unmodified C code on a standard x86-64 processor rather than tuned embedded code.Because a standard PC compiler optimizes code very differently than an embedded compiler, you cannot compare the exact performance numbers—especially the code size—to previous embedded studies \cite{renner,weatherley};. You can only compare the *relative rankings*, meaning which algorithm placed first, second, or third.

\subsection{Ascon Implementations}
Ascon's 5-bit S-box has algebraic degree 2, which makes it unusually cheap to mask against side-channel attacks. Dedicated Ascon hardware and its side-channel behaviour were studied early by Gro\ss{} \emph{et al.} \cite{gross}, and a considerable number of cores have followed, particularly since standardization as NIST SP 800-232 \cite{nist-sp800-232}. The Ascon-AEAD128 core measured in our paper is not intended to compete with these highly optimized, protected versions. Instead, it is a basic, round-per-cycle implementation built specifically to serve as a direct, apples-to-apples control group for our hybrid design using the exact same toolchain.


In summary, this paper presents a controlled substitution rather than a new cipher design. We keep Ascon's mode fixed, replace its standard round function with SIPROUND, and measure the results in C and Verilog against ten existing algorithms using consistent hardware and software flows. The security rationale for the resulting hybrid rests on the maturity of the two components it combines; dedicated cryptanalysis of the combination itself remains future work.